\thispagestyle{plain}

\chapter{Modeling and Evaluation}
\label{cp:mae}

Having prepared our data through rigorous EDA and Feature Engineering, we now move to the training phase. Our objective is to identify
the algorithm that best maximizes the \textbf{Matthews Correlation Coefficient (MCC)}, our chosen metric for this imbalanced dataset.

To ensure a rigorous comparison, we have established a standardized testing protocol that includes advanced hyperparameter optimization,
a robust baseline, and a diverse selection of algorithms ranging from simple linear models to complex ensemble methods.

\section{Optimization Strategy: Why Optuna?}

To unlock the full potential of each algorithm, relying on default hyperparameters is insufficient. While \textit{Grid Search}
is a traditional approach, it suffers from the "curse of dimensionality" and computational inefficiency.

For this project, we used \textbf{Optuna}, a state-of-the-art hyperparameter optimization framework based on \textbf{Bayesian Optimization}.
\begin{itemize}
    \item \textbf{Smart Search:} Unlike Grid Search which tests combinations blindly, Optuna treats optimization as a sequential process.
    It learns from the history of previous trials to construct a probabilistic model, focusing the search on the most promising regions of the
    hyperparameter space.
    \item \textbf{Pruning:} It incorporates an early-stopping mechanism. If a model performs poorly in the early stages of training compared
    to previous trials, it is terminated immediately, freeing up resources to test more configurations.
\end{itemize}
This allows us to fine-tune complex parameters (such as tree depth, learning rates, and regularization) efficiently within a fixed time budget.

\section{Establishing the Floor: The Baseline Model}

Before evaluating predictive models, we must define the "zero-level" performance. We employ a \textbf{Dummy Classifier} using a stratified strategy.
This baseline model generates predictions solely by respecting the class distribution of the training set (randomly predicting $>50K$ $\approx$ 24\%
of the time), completely ignoring the feature data.
\begin{itemize}
    \item \textbf{Goal:} This establishes the minimum MCC score (expected to be $\approx 0$).
    \item \textbf{Implication:} Any valid model must significantly outperform this baseline to prove it has learned actual patterns from the features.
\end{itemize}

\section{Candidate Algorithms}

We have selected a broad spectrum of models to test the linearity and complexity of the dataset. We will present the results for each
category, ranked by their MCC score on the test set.

\subsection{Linear \& Simple Models:}
Designed to test if a simple decision boundary is sufficient.
\begin{itemize}
    \item Logistic Regression (Baseline for linear models)
    \item SGD Classifier (Stochastic Gradient Descent)
    \item SVC (Support Vector Classifier)
    \item KNN (K-Nearest Neighbors)
    \item Gaussian Naive Bayes (Probabilistic approach)
\end{itemize}

\subsection{2. Tree-Based \& Boosting Ensembles:}
Designed to capture non-linearities and complex interactions.
\begin{itemize}
    \item Decision Tree (Single estimator)
    \item Random Forest \& Extra Trees (Bagging)
    \item AdaBoost
    \item \textbf{Gradient Boosting} (Scikit-Learn)
    \item \textbf{XGBoost, LightGBM, CatBoost} (State-of-the-art implementations)
\end{itemize}


\section{Result}

\begin{table}[H]
    \centering
    \label{tab:model_results}
    \begin{tabular}{||l|c|c|c|c||}
        \toprule
        \textbf{Model} & \textbf{Dataset} & \textbf{MCC} & \textbf{AUC-ROC} & \textbf{F1-Weighted} \\
        \midrule
        \midrule
        Dummy Classifier (Baseline) & $\text{X}_{\text{linear}}$ & $ 0.00$ & $ 0.50$ & $ 0.64$ \\
        \midrule
        \textbf{Linear \& Simple Models} & & & & \\
        Logistic Regression & $\text{X}_{\text{linear}}$ & $0.54$ & $0.89$ & $0.83$ \\
        SVC & $\text{X}_{\text{linear}}$ & $0.54$ & $\text{N/A}$ & $0.83$ \\
        Gaussian Naive Bayes & $\text{X}_{\text{linear}}$ & $0.46$ & $0.85$ & $0.78$ \\
        \midrule
        \textbf{Tree-Based Ensembles} & & & & \\
        Random Forest & $\text{X}_{\text{tree}}$ & $0.63$ & $0.92$ & $0.85$ \\
        XGBoost & $\text{X}_{\text{tree}}$ & $ 0.645$ & $0.93$ & $ 0.86$ \\
        LightGBM & $\text{X}_{\text{tree}}$ & $ 0.645$ & $0.93$ & $ 0.86$ \\
        CatBoost & $\text{X}_{\text{tree}}$ & $0.638$ & $0.93$ & $0.86$ \\
        \textbf{Gradient Boosting Classifier} & $\mathbf{X_{\text{tree}}}$ & $\mathbf{0.6477}$ & $\mathbf{0.93}$ & $\mathbf{0.86}$ \\
        \bottomrule
    \end{tabular}
    \caption{Comparison of Model Performance on Test Set}

\end{table}

\clearpage

\section{Stacking}

We built a stacking ensemble using 4 base models and a linear regression meta-learner. Gradient Boosting performed so well on its own that the stacked model quickly reached a performance ceiling around 0.647.

Including too many models in the stack slightly reduced performance, and when we removed the weaker models, performance improved again. This indicates that the least effective models negatively affected the meta-learner instead of contributing useful information.

A few meta-models slightly exceeded this score, but only when using another boosting model as the meta-learner. Since the improvement was minimal compared to the added complexity, we concluded that it was not worth using.

\section{Interpretability}
Performance is not the only goal, understanding the "why" is equally important. Once the best model is identified, we will use
\textbf{Feature Importance analysis} (leveraging XGBoost's gain metrics) to reveal which socioeconomic factors (e.g., Age, Capital, Education)
are the primary drivers of income inequality in this dataset.

